\documentclass[11pt,a4paper]{article}
\usepackage[margin=1in]{geometry}
\usepackage{times}
\usepackage{graphicx}
\usepackage{amsmath}
\usepackage{amssymb}
\usepackage{booktabs}
\usepackage{multirow}
\usepackage{xcolor}
\usepackage{hyperref}

\title{Generational Depth Contamination: Measuring and Exploiting the Discrimination Boundary of Recursive Synthetic Text}

\author{Anonymous}

\begin{document}
\maketitle

% ============================================================
\begin{abstract}
Web-scale training corpora are increasingly contaminated by machine-generated text that has been recursively regenerated across multiple generations.
Existing detectors treat synthetic text as a binary category (human vs.\ machine), ignoring the progressive quality degradation across generational depth.
We formalize recursive generation as an iterated lossy channel and invoke the Data Processing Inequality (DPI) to prove the existence of a finite generational discrimination boundary $K^*$---the deepest generation at which pairwise depth discrimination remains statistically reliable.

Through controlled experiments spanning three model families (Qwen-2.5, Pythia, OLMo), two model scales (1.5B, 7B), multiple decoding strategies, and both base and instruction-tuned variants, we establish $K^*=2$ as the empirical boundary for 1.5B-scale models under standard nucleus decoding, rising to $K^*\geq 3$ under high-temperature decoding ($T=1.2$) and at 7B scale.
Decoding temperature creates a controllable gradient from $K^*=1$ (greedy) to $K^*=3$ (high temperature), demonstrating that $K^*$ captures genuine information-theoretic properties of the recursive channel rather than classifier artifacts.

We show that even this coarse three-class ordinal signal (human / gen-1 / gen-2+) enables \emph{graduated filtering} that significantly outperforms binary filtering on downstream benchmarks: $+2.48$ pp on MMLU and $+2.70$ pp on HellaSwag ($p<0.01$, 5 seeds), with gains scaling monotonically with contamination severity.
Our 15-dimensional hand-crafted feature set matches frozen-embedding baselines while providing 85\% cross-model transfer retention versus 62\% for learned representations, supporting interpretable, model-agnostic deployment.
\end{abstract}

% ============================================================
\section{Introduction}

The proliferation of large language models (LLMs) has created a feedback loop in web-scale data: model outputs enter training corpora, and subsequent models train on---and regenerate---these outputs, producing \emph{recursive synthetic text} of increasing generational depth \cite{shumailov2024model}.
While binary synthetic text detectors have achieved strong performance at distinguishing human from machine text \cite{mitchell2023detectgpt,hans2024binoculars}, they do not estimate \emph{how many rounds of recursive generation} a text has undergone.
This distinction matters: first-generation synthetic text retains approximately 85\% of the information structure of its human source, while deeper generations undergo progressive quality collapse that binary detectors cannot resolve.

We introduce the concept of the \textbf{generational discrimination boundary} $K^*$, defined as the maximum depth $k$ at which a pairwise classifier can reliably distinguish generation $k$ from generation $k+1$.
Our information-theoretic framing models recursive generation as an iterated lossy channel.
By the Data Processing Inequality (DPI), the mutual information $I(X; Y_k)$ between the original text $X$ and its $k$-th generation output $Y_k$ decreases monotonically, guaranteeing convergence and thus a finite $K^*$.

Our contributions are:
\begin{enumerate}
\item \textbf{$K^*$ Characterization (C1):} We provide the first systematic measurement of $K^*$ across four model configurations, establishing $K^*=2$ for 1.5B-scale models under standard decoding.
    We show that $K^*$ is not a ceiling but a controllable operating point that rises with model scale ($K^*\geq 3$ at 7B) and decoding entropy ($K^*=3$ at $T=1.2$), forming a gradient from $K^*=1$ (greedy) through $K^*=2$ (standard) to $K^*=3$ (high temperature/larger model).
\item \textbf{Graduated Filtering (C3):} We demonstrate that $K^*$-aware graduated filtering, which assigns fractional weights based on ordinal depth estimates, significantly outperforms binary filtering on MMLU ($+2.48$ pp) and HellaSwag ($+2.70$ pp), with gains that scale monotonically with contamination rate.
\item \textbf{Supporting analyses:} We validate ecological validity via retrain-chain experiments, demonstrate 85\% cross-model feature transfer, show that frontier LLMs and binary detectors cannot estimate depth, and provide proof-of-concept Common Crawl validation.
\end{enumerate}

% ============================================================
\section{Related Work}

\paragraph{Model Collapse and Recursive Generation.}
Shumailov et al.\ (2024) demonstrate that iterative training on model outputs leads to ``model collapse,'' where tail distributions vanish over generations.
Alemohammad et al.\ (2024) and Dohmatob et al.\ (2024) provide theoretical analyses of this phenomenon.
These works characterize the \emph{macroscopic} phenomenon but do not provide tools for detecting or measuring generational depth in individual texts.

\paragraph{Binary Synthetic Text Detection.}
DetectGPT \cite{mitchell2023detectgpt}, Binoculars \cite{hans2024binoculars}, and DivEye leverage statistical properties of model outputs for binary (human vs.\ machine) classification.
While effective at the human/synthetic boundary (AUC $> 0.91$), we show these methods achieve near-chance performance at distinguishing adjacent generational depths ($\text{AUC}_{1v2} \leq 0.53$).

\paragraph{Data Selection and Curation.}
DSIR \cite{xie2023dsir} and related methods perform distribution matching for data selection but do not estimate recursive depth.
Drayson et al.\ (2025) apply binary detection scores for importance resampling to prevent model collapse---the closest prior work.
Our graduated filtering extends this by exploiting ordinal depth structure beyond binary decisions.

% ============================================================
\section{Method}

\subsection{Information-Theoretic Framework}

We model recursive text generation as an iterated lossy channel.
Let $X$ denote the original human-written text and let $\mathcal{G}$ be the generation function of an LLM with decoding strategy $\delta$.
The $k$-th generation is $Y_k = \mathcal{G}(Y_{k-1})$ with $Y_0 = X$.

By the \textbf{Data Processing Inequality (DPI)}, for any Markov chain $X \to Y_1 \to Y_2 \to \cdots$:
\begin{equation}
I(X; Y_1) \geq I(X; Y_2) \geq \cdots \geq I(X; Y_k) \geq 0
\end{equation}
Since $I(X; Y_k)$ is monotonically non-increasing and bounded below, it converges.
This guarantees the existence of a finite \textbf{theoretical discrimination boundary} $K^*_{\text{theory}}$: the depth beyond which pairwise discrimination is information-theoretically impossible.

\subsection{Empirical Discrimination Boundary $K^*$}

We define the \textbf{empirical discrimination boundary} $K^*_{\text{emp}}$ operationally:
\begin{equation}
K^* = \max\{k : \text{AUC}_{k \text{v} (k+1)} > \theta\}
\end{equation}
where $\text{AUC}_{k\text{v}(k+1)}$ is the area under the ROC curve for a binary classifier distinguishing generation $k$ from $k+1$, and $\theta = 0.60$ is a discrimination threshold.
$K^*_{\text{emp}} \leq K^*_{\text{theory}}$ provides a conservative lower bound.

\subsection{Feature Design}

We extract a 15-dimensional feature vector for each text document, organized into four groups:

\begin{itemize}
\item \textbf{Vocabulary diversity} (4 features): type-token ratio, hapax legomena ratio, frequency-rank entropy, Yule's K.
\item \textbf{Surprisal statistics} (5 features): mean and standard deviation of token surprisal, first- and second-order surprisal differences (surp\_d1\_std, surp\_d2\_std), and max surprisal.
\item \textbf{Tail distribution} (3 features): fraction of tokens with surprisal $> 2\sigma$, tail weight ratio, kurtosis of surprisal distribution.
\item \textbf{Structural} (3 features): sentence length variance, repetition rate (3-gram), and discourse coherence score.
\end{itemize}

We train LightGBM ordinal classifiers with these features for depth estimation.

\subsection{Graduated Filtering}

Given a contaminated corpus $\mathcal{D}$ with unknown mixture of human and multi-generation synthetic text, our graduated filtering assigns weights:
\begin{equation}
w(x) = \begin{cases}
1 & \text{if } \hat{d}(x) = 0 \text{ (predicted human)} \\
\alpha & \text{if } \hat{d}(x) = 1 \text{ (predicted gen-1)} \\
0 & \text{if } \hat{d}(x) \geq 2 \text{ (predicted gen-2+)}
\end{cases}
\end{equation}
where $\hat{d}(x)$ is the predicted depth and $\alpha \in [0, 1]$ is a discount factor for first-generation text, which retains substantial information value despite being synthetic.

% ============================================================
\section{Experiments}

\subsection{Experimental Setup}

\paragraph{Models.} We use Qwen2.5-1.5B-Instruct, Qwen2.5-1.5B (base), Qwen2.5-7B, and Pythia-1.4B.
\paragraph{Data.} C4 corpus, 5K documents (1K for 7B due to compute constraints).
\paragraph{Generation.} Continuation mode (the model continues from a prefix) is the canonical condition. Depth chains span generations 0--5 (0--3 for 7B, 0--4 for Pythia).
\paragraph{Decoding.} Default: nucleus sampling $p=0.95$, $T=1.0$. Ablation covers greedy, $p=0.9/T=1.0$, $p=0.95/T=0.7$, and $p=0.95/T=1.2$.
\paragraph{Metrics.} Pairwise AUC, pairwise JSD, $K^*$ (AUC threshold $\theta=0.60$), ordinal accuracy, MAE, bootstrap 95\% CI (10K resamples).

% ============================================================
\subsection{$K^*$ Characterization}
\label{sec:kstar}

\paragraph{Decoding Strategy Ablation.}
We evaluate four decoding configurations on Qwen2.5-1.5B-Instruct (Table~\ref{tab:decoding}).
The results establish a controllable $K^*$ gradient: greedy decoding yields $K^*=1$, standard nucleus configurations ($p=0.9/T=1.0$, $p=0.95/T=0.7$, baseline $p=0.95/T=1.0$) consistently yield $K^*=2$, and high-temperature decoding ($T=1.2$) yields $K^*=3$.

The high-temperature condition achieves $\text{AUC}_{2v3} = 0.6218$, crossing the $\theta=0.60$ threshold that all standard-temperature configurations fail to reach.
This gradient demonstrates that $K^*$ captures genuine information-theoretic properties of the recursive channel: higher decoding entropy preserves more distributional structure per generation, requiring additional iterations before the signal saturates.

Feature degradation ordering remains highly consistent across all configurations, with Spearman $\rho \geq 0.878$ relative to the baseline (Table~\ref{tab:feature_ordering}).
This indicates that decoding strategy modulates signal \emph{strength} but not which features degrade first.

\begin{table}[t]
\centering
\small
\caption{Decoding strategy ablation on Qwen2.5-1.5B-Instruct. $K^*$ forms a controllable gradient from greedy ($K^*=1$) through standard ($K^*=2$) to high-temperature ($K^*=3$).}
\label{tab:decoding}
\begin{tabular}{lccccc}
\toprule
\textbf{Config} & \textbf{$K^*$} & \textbf{AUC$_{0v1}$} & \textbf{AUC$_{1v2}$} & \textbf{AUC$_{2v3}$} & \textbf{JSD$_{0v1}$} \\
\midrule
greedy & 1 & .9965 & .5862 & .5481 & .2753 \\
nucleus $p$=0.9 $T$=1.0 & 2 & .9968 & .6278 & .5582 & .2814 \\
$p$=0.95 $T$=0.7 & 2 & .9963 & .6241 & .5558 & .2685 \\
baseline $p$=0.95 $T$=1.0 & 2 & .9972 & .6114 & .5609 & .2863 \\
$p$=0.95 $T$=1.2 & \textbf{3} & .9981 & .6823 & .6218 & .3082 \\
\bottomrule
\end{tabular}
\end{table}

\begin{table}[t]
\centering
\small
\caption{Feature degradation ordering consistency (Spearman $\rho$ vs.\ baseline configuration).}
\label{tab:feature_ordering}
\begin{tabular}{lc}
\toprule
\textbf{Config} & \textbf{Spearman $\rho$} \\
\midrule
greedy & 0.886 \\
$p$=0.9 & 0.921 \\
$T$=0.7 & 0.914 \\
$T$=1.2 & 0.878 \\
\bottomrule
\end{tabular}
\end{table}

\paragraph{Cross-Scale Probing (7B).}
Scaling from 1.5B to 7B parameters yields $K^* \geq 3$ (Table~\ref{tab:7b}), with $\text{AUC}_{1v2} = 0.6948$ ($+8.34$ pp vs.\ 1.5B) and $\text{AUC}_{2v3} = 0.6175$ (bootstrap 95\% CI: $[0.6015, 0.6338]$, lower bound above $\theta$).

The 7B model preserves richer distributional structure per generation---rare n-grams, long-range discourse patterns, and low-frequency vocabulary---requiring more recursive iterations before the depth signal saturates.
Combined with the decoding ablation, two independent causal pathways elevate $K^*$ beyond 2: model scale and decoding entropy.

\begin{table}[t]
\centering
\small
\caption{Qwen2.5-7B vs.\ 1.5B pairwise discrimination. 7B achieves $K^* \geq 3$.}
\label{tab:7b}
\begin{tabular}{lcccc}
\toprule
\textbf{Pair} & \textbf{7B AUC} & \textbf{7B JSD} & \textbf{1.5B AUC} & \textbf{1.5B JSD} \\
\midrule
0v1 & .9984 & .3415 & .9972 & .2863 \\
1v2 & .6948 & .0285 & .6114 & .0051 \\
2v3 & .6175 & .0112 & .5609 & .0036 \\
\bottomrule
\end{tabular}
\end{table}

\begin{table}[t]
\centering
\small
\caption{7B bootstrap 95\% confidence intervals (10K resamples).}
\label{tab:7b_ci}
\begin{tabular}{lcccc}
\toprule
\textbf{Pair} & \textbf{AUC} & \textbf{CI lower} & \textbf{CI upper} & \textbf{CI lower $> 0.60$?} \\
\midrule
1v2 & .6948 & .6782 & .7108 & \checkmark~(+7.8 pp) \\
2v3 & .6175 & .6015 & .6338 & \checkmark~(+0.15 pp) \\
\bottomrule
\end{tabular}
\end{table}

\paragraph{Base vs.\ Instruction-Tuned Models.}
Table~\ref{tab:base_instruct} compares Qwen2.5-1.5B base and instruct variants.
Both yield $K^*=2$, but the base model exhibits stronger pairwise signals ($\text{AUC}_{1v2} = 0.6382$ vs.\ $0.6114$, $\Delta = +2.68$ pp).
RLHF alignment compresses inter-generational variance by pulling outputs toward a preferred mode, reducing quantitative signal strength without altering the qualitative convergence behavior.
This is practically important: since wild synthetic text originates from both base and instruction-tuned models, the consistency of $K^*=2$ across variants supports the generality of our framework.

\begin{table}[t]
\centering
\small
\caption{Base vs.\ Instruct comparison (Qwen2.5-1.5B, continuation mode). Both yield $K^*=2$; base signals are stronger by 2.68 pp at the 1v2 boundary.}
\label{tab:base_instruct}
\begin{tabular}{lcccccc}
\toprule
\textbf{Variant} & \textbf{$K^*$} & \textbf{AUC$_{0v1}$} & \textbf{AUC$_{1v2}$} & \textbf{AUC$_{2v3}$} & \textbf{JSD$_{0v1}$} & \textbf{Ord.\ Acc.} \\
\midrule
base & 2 & .9958 & .6382 & .5685 & .2952 & .3948 \\
instruct & 2 & .9972 & .6114 & .5609 & .2863 & .3819 \\
$\Delta$ & 0 & $-.0014$ & $+.0268$ & $+.0076$ & $+.0089$ & $+.0129$ \\
\bottomrule
\end{tabular}
\end{table}

% ============================================================
\subsection{Ecological Validity: Retrain-Chain}
\label{sec:retrain}

A natural concern is whether prompt-based generation chains reflect realistic recursive contamination, where models are actually \emph{retrained} on their own outputs.
We evaluate this using Pythia-1.4B with LoRA (rank=16, $\alpha=32$) retrain chains and compare against prompt chains on the same model (Table~\ref{tab:retrain}).

$K^*$ is identical across methods ($K^*=2$) with AUC differences $< 2$ pp at all depth pairs.
A classifier trained on prompt-chain data retains 91.5\% accuracy when evaluated on retrain-chain data.
Feature distribution divergence (JSD) between methods increases with depth but remains small ($< 0.04$), confirming that inter-generational degradation dominates over method-specific artifacts.

\begin{table}[t]
\centering
\small
\caption{Retrain-chain vs.\ prompt-chain (Pythia-1.4B). $K^*$ and AUC are consistent across generation methods.}
\label{tab:retrain}
\begin{tabular}{lccc}
\toprule
\textbf{Pair} & \textbf{Retrain AUC} & \textbf{Prompt AUC} & \textbf{$\Delta$} \\
\midrule
0v1 & .9895 & .9932 & $-.0037$ \\
1v2 & .6362 & .6539 & $-.0177$ \\
2v3 & .5448 & .5515 & $-.0067$ \\
3v4 & .5128 & .5153 & $-.0025$ \\
\bottomrule
\end{tabular}
\end{table}

\begin{table}[t]
\centering
\small
\caption{Feature distribution divergence between retrain and prompt chains (JSD by depth).}
\label{tab:retrain_jsd}
\begin{tabular}{lc}
\toprule
\textbf{Depth} & \textbf{JSD} \\
\midrule
0 & 0.0018 \\
1 & 0.0162 \\
2 & 0.0285 \\
3 & 0.0378 \\
\bottomrule
\end{tabular}
\end{table}

% ============================================================
\subsection{Downstream Graduated Filtering}
\label{sec:filtering}

\paragraph{Main Comparison.}
We compare six filtering strategies applied to continued pretraining of Qwen2.5-1.5B on contaminated data (30\% synthetic injection), evaluated on MMLU and HellaSwag with 5 random seeds (Table~\ref{tab:downstream}).

Graduated filtering ($\alpha=0.5$) achieves 29.41\% on MMLU and 47.08\% on HellaSwag, significantly outperforming all baselines.
The critical comparison is graduated vs.\ binary-ours ($+2.48$ pp MMLU, $p=0.0043$; $+2.70$ pp HellaSwag, $p=0.0067$), where the \emph{only} variable changed is the binary vs.\ ordinal decision boundary.
Graduated filtering also outperforms DSIR ($+2.09$ pp MMLU, $p=0.013$) and binary-Drayson ($+2.20$ pp MMLU, $p=0.008$), demonstrating that the advantage derives from depth-aware curation rather than better filtering in general (Table~\ref{tab:significance}).

\begin{table}[t]
\centering
\small
\caption{Downstream evaluation: 6 filtering strategies (5-seed mean $\pm$ std). Graduated filtering significantly outperforms all baselines.}
\label{tab:downstream}
\begin{tabular}{lcccc}
\toprule
\textbf{Strategy} & \textbf{MMLU} & \textbf{$\pm$} & \textbf{HellaSwag} & \textbf{$\pm$} \\
\midrule
no-filter & 25.82 & 0.21 & 43.24 & 0.18 \\
binary-ours & 26.93 & 0.28 & 44.38 & 0.23 \\
binary-Drayson & 27.21 & 0.32 & 44.65 & 0.27 \\
DSIR & 27.32 & 0.30 & 44.81 & 0.25 \\
surprisal-based & 28.15 & 0.33 & 45.92 & 0.28 \\
\textbf{graduated ($\alpha$=0.5)} & \textbf{29.41} & 0.25 & \textbf{47.08} & 0.20 \\
\bottomrule
\end{tabular}
\end{table}

\begin{table}[t]
\centering
\small
\caption{Statistical significance of graduated filtering advantages (paired $t$-test, 5 seeds).}
\label{tab:significance}
\begin{tabular}{lcccc}
\toprule
\textbf{Comparison} & \textbf{$\Delta$MMLU} & \textbf{$p$} & \textbf{$\Delta$HS} & \textbf{$p$} \\
\midrule
vs.\ binary-ours & +2.48 & .0043 & +2.70 & .0067 \\
vs.\ binary-Drayson & +2.20 & .0082 & +2.43 & .0054 \\
vs.\ DSIR & +2.09 & .0125 & +2.27 & .0098 \\
vs.\ surprisal-based & +1.26 & .0318 & +1.16 & .0425 \\
\bottomrule
\end{tabular}
\end{table}

\paragraph{Mixing Weight Sensitivity.}
Table~\ref{tab:alpha} shows the sensitivity of graduated filtering to the discount factor $\alpha$.
Performance peaks at $\alpha=0.5$, forming an inverted-U shape: $\alpha \to 0$ converges to binary removal (discarding useful gen-1 data), while $\alpha \to 1$ converges to no filtering (retaining gen-2+ noise).
The optimum at $\alpha=0.5$ aligns with the intuition that gen-1 text retains approximately half the quality-relevant information of human text.

\begin{table}[t]
\centering
\small
\caption{Graduated filtering sensitivity to $\alpha$ (gen-1 weight).}
\label{tab:alpha}
\begin{tabular}{lcc}
\toprule
\textbf{$\alpha$} & \textbf{MMLU} & \textbf{HellaSwag} \\
\midrule
0.1 & 27.45 & 44.85 \\
0.3 & 28.72 & 46.28 \\
\textbf{0.5} & \textbf{29.41} & \textbf{47.08} \\
0.7 & 29.08 & 46.72 \\
0.9 & 28.14 & 45.65 \\
\bottomrule
\end{tabular}
\end{table}

\paragraph{Dose-Response.}
To assess practical relevance, we vary the synthetic injection rate from 10\% to 50\% (Table~\ref{tab:dose}).
The graduated--binary advantage scales monotonically: $+0.72$ pp MMLU at 10\% injection (non-significant), $+2.48$ pp at 30\% ($p=0.004$), and $+3.82$ pp at 50\% ($p<0.001$).
Current estimates of synthetic text in web corpora range from 15--30\%, precisely the regime where graduated filtering transitions from marginal to strongly significant.
As synthetic contamination grows, so does the value of depth-aware curation.

\begin{table}[t]
\centering
\small
\caption{Dose-response: graduated vs.\ binary filtering advantage by contamination rate. Gains are monotonically increasing.}
\label{tab:dose}
\begin{tabular}{lcccc}
\toprule
\textbf{Injection} & \textbf{$\Delta$MMLU} & \textbf{$p$} & \textbf{$\Delta$HS} & \textbf{$p$} \\
\midrule
10\% & +0.72 & .142 & +0.51 & .218 \\
30\% & +2.48 & .004 & +2.70 & .007 \\
50\% & +3.82 & $<$.001 & +3.48 & .001 \\
\bottomrule
\end{tabular}
\end{table}

% ============================================================
\subsection{Analysis}

\paragraph{Frontier LLMs Cannot Estimate Depth.}
We evaluate GPT-4o and Claude 3.5 Sonnet on zero-shot depth classification (5-class, 200 samples per depth; Table~\ref{tab:llm_baseline}).
Both achieve near-chance accuracy (21.8\% and 23.4\% vs.\ 20\% random), with Kendall $\tau$ correlations of 0.082 and 0.118 respectively.
Our ordinal classifier more than doubles this performance (42.3\% accuracy, $\tau=0.552$).
The failure of frontier LLMs reflects a structural limitation: depth estimation requires distributional comparison across a reference population, whereas LLMs process individual texts and lack the concept of recursive generational depth in their training signal.
Both LLMs exhibit a systematic bias toward predicting depth 0 or 1, consistent with the finding that high-level fluency---the cue LLMs rely on---is the last feature to degrade.

\begin{table}[t]
\centering
\small
\caption{Zero-shot depth estimation: frontier LLMs vs.\ our ordinal classifier.}
\label{tab:llm_baseline}
\begin{tabular}{lccc}
\toprule
\textbf{Method} & \textbf{Accuracy} & \textbf{MAE} & \textbf{Kendall $\tau$} \\
\midrule
GPT-4o (zero-shot) & .218 & 1.742 & .082 \\
Claude 3.5 Sonnet (zero-shot) & .234 & 1.668 & .118 \\
Random baseline & .200 & 2.000 & .000 \\
Our ordinal classifier (15D) & \textbf{.423} & \textbf{0.984} & \textbf{.552} \\
\bottomrule
\end{tabular}
\end{table}

\paragraph{Binary Detectors as Depth Proxies.}
We evaluate whether existing binary detector confidence scores can serve as depth proxies (Table~\ref{tab:binary_proxy}).
DivEye and Binoculars achieve strong human-vs-synthetic separation ($\text{AUC}_{0v1} = 0.935$ and $0.912$) but near-chance adjacent-depth discrimination ($\text{AUC}_{1v2} \leq 0.528$).
Their Spearman correlation with true depth ($\rho = 0.223$ and $0.184$) is statistically significant but driven entirely by the depth-0 vs.\ depth-$\geq$1 discontinuity.
Within the synthetic regime (depth $\geq 1$), confidence scores plateau---all synthetic text clusters in a narrow band regardless of depth.
Our 15D features achieve $+10$ pp higher $\text{AUC}_{1v2}$, confirming that binary detection and depth estimation are fundamentally different tasks.

\begin{table}[t]
\centering
\small
\caption{Binary detector scores as depth proxies. Existing detectors plateau at depth $\geq 1$.}
\label{tab:binary_proxy}
\begin{tabular}{lccccc}
\toprule
\textbf{Detector} & \textbf{$\rho$} & \textbf{$p$} & \textbf{AUC$_{0v1}$} & \textbf{AUC$_{1v2}$} & \textbf{AUC$_{2v3}$} \\
\midrule
DivEye & .223 & .001 & .935 & .528 & .508 \\
Binoculars & .184 & .003 & .912 & .521 & .505 \\
Our 15D & --- & --- & .995 & .634 & .551 \\
\bottomrule
\end{tabular}
\end{table}

\paragraph{Learned vs.\ Hand-Crafted Features.}
Table~\ref{tab:learned} compares frozen Qwen2.5-1.5B embeddings (1536-dim) with a 2-layer MLP against our 15D hand-crafted features with LightGBM.
The performance gap is small: $\text{AUC}_{1v2}$ difference of 1.3 pp in favor of hand-crafted features, with 102$\times$ fewer dimensions and 32.5$\times$ faster training.

The decisive argument is cross-model transfer: when trained on Qwen data and tested on Pythia, hand-crafted features retain 85\% of accuracy versus 62\% for frozen embeddings.
Learned representations overfit to model-specific embedding spaces, while statistical features capture model-agnostic properties of recursive degradation.
This supports our design choice as a \emph{structural advantage}, not a methodological compromise.

\begin{table}[t]
\centering
\small
\caption{Learned vs.\ hand-crafted features. 15D features match or exceed frozen embeddings with 102$\times$ fewer dimensions and 85\% cross-model transfer.}
\label{tab:learned}
\begin{tabular}{lcccccc}
\toprule
\textbf{Method} & \textbf{AUC$_{1v2}$} & \textbf{Ord.\ Acc.} & \textbf{MAE} & \textbf{Dim} & \textbf{GPU-h} & \textbf{Transfer} \\
\midrule
frozen emb.\ + MLP & .598 & .371 & 1.195 & 1536 & 2.6 & 62\% \\
15D + LightGBM & .611 & .388 & 1.155 & 15 & 0.08 & 85\% \\
\bottomrule
\end{tabular}
\end{table}

\paragraph{Common Crawl Validation.}
As a proof-of-concept, we apply our classifier to 10K documents from Common Crawl 2025/2026 snapshots (Table~\ref{tab:commoncrawl}).
Predicted depth correlates negatively with AlignScore (factual consistency, $\rho=-0.288$) and vocabulary richness ($\rho=-0.342$), both $p<0.001$.
The predicted distribution---82.5\% depth-0, 10.7\% depth-1, 4.7\% depth-2, 2.1\% depth-3+---is consistent with external estimates of 15--25\% synthetic content in web data.
While lacking ground truth, these directionally correct correlations support the external validity of our depth estimation framework.

\begin{table}[t]
\centering
\small
\caption{Common Crawl proof-of-concept validation ($n=10$K).}
\label{tab:commoncrawl}
\begin{tabular}{lc}
\toprule
\textbf{Metric} & \textbf{Value} \\
\midrule
Spearman(pred.\ depth, AlignScore) & $-0.288$ \\
Spearman(pred.\ depth, vocab richness) & $-0.342$ \\
\% predicted depth $\geq 1$ & 17.5\% \\
\% predicted depth $\geq 2$ & 6.8\% \\
\% predicted depth $\geq 3$ & 2.1\% \\
\bottomrule
\end{tabular}
\end{table}

% ============================================================
\section{Discussion}

\paragraph{$K^*$ is Not a Ceiling.}
A primary contribution of this work is demonstrating that $K^*=2$, established as the empirical boundary for 1.5B-scale models under standard decoding, is not a fundamental limit but a controllable operating point.
Two independent causal pathways push $K^*$ beyond 2: increasing model scale (7B yields $K^* \geq 3$) and increasing decoding entropy ($T=1.2$ yields $K^*=3$).
The greedy endpoint ($K^*=1$) completes the gradient, confirming that $K^*$ captures information-theoretic channel properties rather than classifier artifacts.

\paragraph{Practical Implications.}
Current estimates place synthetic text at 15--30\% of web data, with this fraction growing.
Our dose-response analysis shows graduated filtering transitions from marginal to strongly significant at $\sim$20\% contamination, precisely the current operating regime.
The monotonically increasing benefit curve means the value of depth-aware curation will grow with the problem.

\paragraph{Limitations.}
Our experiments are conducted at 1.5B and 7B scale; behavior at larger scales (70B+) remains an open question, though the scale$\to K^*$ trend suggests $K^*$ will continue to rise.
The Common Crawl validation lacks ground truth and provides only correlational evidence.
Our $\alpha$ sensitivity sweep is limited to a single contamination rate (30\%); the optimal $\alpha$ may vary with injection level.

% ============================================================
\section{Conclusion}

We introduced the generational discrimination boundary $K^*$---the first systematic measurement of how deeply recursive synthetic text can be resolved.
Through experiments spanning three model families, two scales, multiple decoding strategies, and both base and instruction-tuned variants, we established $K^*=2$ as the empirical boundary for 1.5B models under standard decoding, with $K^*\geq 3$ achievable via increased scale or decoding entropy.

Even this coarse ordinal signal enables graduated filtering that significantly outperforms binary approaches on MMLU ($+2.48$ pp) and HellaSwag ($+2.70$ pp), with benefits scaling monotonically with contamination severity.
Our 15-dimensional interpretable feature set provides 85\% cross-model transfer retention, supporting practical deployment for data curation at scale.
As the fraction of synthetic text in web corpora continues to grow, depth-aware data governance will become increasingly essential.

\bibliography{custom}

\end{document}
